% =====================================================================
%  PART I -- THE STATE SPACE
% =====================================================================

\chapter{The Base Complex}\label{ch:complex}

\section{What the object is}

\begin{intu}\Intu
Picture a research assistant that is not one mind but a small society of
specialists: one that searches, one that remembers, one that reasons, one
that checks. \textsc{tatiana} is the mathematics of \emph{that society} ---
who is bound to whom, whether their beliefs agree, by how much they
disagree and exactly where, and how the society rewires itself as it
works. The reasoning itself is done by an external language model. The
mathematics governs the society, not the neurons.
\end{intu}

\begin{definition}[Cognitive state space]\label{def:M}
A \emph{cognitive state space} is a triple $\M=(\CC,\Fs,s)$ where
\begin{enumerate}[label=(\roman*)]
  \item $\CC$ is a stratified finite abstract simplicial complex
        (Section~\ref{sec:strat});
  \item $\Fs$ is a cellular sheaf of finite-dimensional real
        inner-product spaces on $\CC$ (Chapter~\ref{ch:sheaf});
  \item $s\in\Czero$ is a distinguished $0$-cochain, the \emph{current
        state}, whose coherence is measured by the sheaf Laplacian
        (Chapter~\ref{ch:laplacian}).
\end{enumerate}
The space evolves in discrete time under structural operators drawn from
an operator algebra $\Dop$ which is itself permitted to grow
(Part~V).
\end{definition}

\source{\texttt{DOCS/TATIANA\_Mathematical\_Description.tex}, Def.~1.1,
merged with \texttt{DOCS/TATIANA\_LOGBOOK.md} \S0.1. The logbook writes
the triple as $(C,F,s)$ with $C$ ``a stratified simplicial complex
(coarse level = cognitive organs as vertices; fine level = concepts
inside each organ's stalk)''. Note that this phrasing conflates two
different things --- the fine complex $\CC_v$ and the stalk $\stalk{v}$
--- which are kept strictly separate here and in
Chapter~\ref{ch:sheaf}. See Appendix~\ref{app:discrepancy}, item D1.}

\section{Simplicial complexes and $n$-ary binding}

\begin{intu}\Intu
A graph records \emph{pairwise} facts: $A$ relates to $B$. But a theorem
depending on three hypotheses is \emph{one} three-way fact, not three
separate pairs --- drop any one hypothesis and the fact is gone, which is
not how three independent edges behave. A simplicial complex is the
smallest honest language for facts of arbitrary arity: a filled triangle
is a genuine three-way binding.
\end{intu}

\begin{definition}[Abstract simplicial complex]\label{def:asc}
Let $V$ be a finite set. An \emph{abstract simplicial complex} $\CC$ on
$V$ is a family of nonempty subsets of $V$, called \emph{simplices},
which is \emph{downward closed}: if $\sigma\in\CC$ and
$\emptyset\neq\tau\subseteq\sigma$, then $\tau\in\CC$. A simplex of
cardinality $k+1$ is a \emph{$k$-simplex}; write $\dim\sigma=k$. If
$\tau\subseteq\sigma$ we write $\tau\triang\sigma$ and call $\tau$ a
\emph{face} of $\sigma$. The set of $k$-simplices is denoted $\CC_k$.
\end{definition}

Thus $0$-simplices are vertices, $1$-simplices edges, $2$-simplices
filled triangles, and so on. Downward closure is not a technical
convenience; it is the modelling assumption that every sub-collection of
a jointly bound fact is itself recorded.

\begin{remark}[Simplices are $n$-ary relations]\label{rem:nary}
A $k$-simplex $\{v_0,\dots,v_k\}$ models a single $(k{+}1)$-ary binding
among its vertices, not $\binom{k+1}{2}$ pairwise ones. This distinction
is what later earns the use of $b_1$ as a diagnostic
(Part~IV): a $1$-cycle counts as a genuine structural hole
only because a \emph{filled} $2$-simplex means ``these three are jointly
bound'' rather than merely ``each pair is bound''.
\end{remark}

\begin{remark}[Two relations on the same vertices]\label{rem:twostructures}
A logical prerequisite chain is \emph{directed}: $A$ is needed for $B$.
A simplicial complex is undirected. \textsc{tatiana} therefore carries
two relations on the same vertex set --- a directed poset for logical
dependency, and the undirected simplices for $n$-ary co-binding. They are
never conflated, and no result in this book uses one where the other is
meant.
\end{remark}

\section{The two-level stratification}\label{sec:strat}

\begin{intu}\Intu
The society has two scales. Zoom out and the vertices are whole
\emph{organs}. Zoom into one organ and it is itself a little complex whose
vertices are \emph{concepts}. The same geometry runs at both scales, and
the two are glued by a coarse-graining map developed in Part~III.
\end{intu}

\begin{construction}[The stratified complex]\label{con:strat}
$\CC$ carries two strata.
\begin{itemize}
  \item The \emph{coarse complex} $\Kc$ has as vertices the cognitive
        organs
        \[
          V(\Kc)=\{\textsc{Planner},\ \textsc{Search},\
          \textsc{Canonicalizer},\ \textsc{Curator},\
          \textsc{WorkingMemory},\ \textsc{Reasoning},\ \textsc{Verify}\}.
        \]
        A simplex $\sigma\in\Kc$ is a \emph{bound coalition} of organs
        acting as a unit, carrying a time-dependent coupling weight
        $w(\sigma,t)\in[0,1]$.
  \item For each organ $v\in V(\Kc)$, a \emph{fine complex} $\CC_v$ has
        concepts as vertices and $n$-ary conceptual bindings as
        simplices. This is the organ's internal knowledge.
\end{itemize}
\end{construction}

The stratification answers an ambiguity present in the earliest drafts
--- \emph{are the vertices concepts or subsystems?} --- by answering
\emph{both}, at two different levels.

\begin{hon}\Hon
\Contested\ The word ``stratified'' is used here in a weak,
combinatorial sense: two nested scales with a map between them. It does
\emph{not} refer to a Whitney stratification, a stratified space in the
sense of intersection homology, or a constructible sheaf on a stratified
base, and no theorem from those theories is invoked anywhere in this
book. Earlier drafts borrowed vocabulary from that literature; the
logbook (\S0.5) records that most such borrowing was cut as decorative.
This is a naming collision, and the reader should not import
expectations from it.
\end{hon}

\section{The running example}\label{sec:running}

Every instrument built in this book will be applied to one small
complex, so that the same object is measured repeatedly and the
measurements can be compared.

\begin{example}[The complex $\Kc_4$]\label{ex:K4}
Let $\Kc_4$ have vertex set
\[
  V=\{\mathsf{S},\ \mathsf{M},\ \mathsf{R},\ \mathsf{V}\}
\]
(read: \textsc{Search}, \textsc{Memory}, \textsc{Reasoning},
\textsc{Verify}), with $1$-simplices
\[
  E=\bigl\{\{\mathsf{S},\mathsf{M}\},\ \{\mathsf{M},\mathsf{R}\},\
  \{\mathsf{S},\mathsf{R}\},\ \{\mathsf{R},\mathsf{V}\}\bigr\}
\]
and a single $2$-simplex $\{\mathsf{S},\mathsf{M},\mathsf{R}\}$, filled.
Downward closure holds: all three faces of the triangle are in $E$.
\end{example}

\begin{center}
\begin{tikzpicture}[scale=1.6,
  vtx/.style={circle,fill=black,inner sep=1.8pt},
  lbl/.style={font=\small}]
  \coordinate (S) at (0,1);
  \coordinate (M) at (-1,0);
  \coordinate (R) at (1,0);
  \coordinate (V) at (2.3,-0.6);
  \fill[linkblue,opacity=0.18] (S) -- (M) -- (R) -- cycle;
  \draw[thick] (S)--(M) (M)--(R) (S)--(R) (R)--(V);
  \node[vtx] at (S) {}; \node[vtx] at (M) {};
  \node[vtx] at (R) {}; \node[vtx] at (V) {};
  \node[lbl,above] at (S) {$\mathsf{S}$};
  \node[lbl,left]  at (M) {$\mathsf{M}$};
  \node[lbl,below right] at (R) {$\mathsf{R}$};
  \node[lbl,right] at (V) {$\mathsf{V}$};
  \node[font=\footnotesize,linkblue] at (0,0.35) {filled};
\end{tikzpicture}
\captionof{figure}{The running example $\Kc_4$. The triangle
$\{\mathsf{S},\mathsf{M},\mathsf{R}\}$ is filled: it is one ternary
binding, not three binary ones. The edge to $\mathsf{V}$ is a pendant.}
\label{fig:K4}
\end{center}

Its degrees are
$\deg\mathsf{S}=2$, $\deg\mathsf{M}=2$, $\deg\mathsf{R}=3$,
$\deg\mathsf{V}=1$, and its graph Laplacian, in the vertex order
$(\mathsf{S},\mathsf{M},\mathsf{R},\mathsf{V})$, is
\begin{equation}\label{eq:LK4}
  L_{\Kc_4}=
  \begin{pmatrix}
     2 & -1 & -1 &  0\\
    -1 &  2 & -1 &  0\\
    -1 & -1 &  3 & -1\\
     0 &  0 & -1 &  1
  \end{pmatrix}.
\end{equation}

\begin{proposition}[Spectrum of the running example]\label{prop:specK4}
\Proved\ The characteristic polynomial of $L_{\Kc_4}$ is
$\lambda(\lambda-1)(\lambda-3)(\lambda-4)$, so
\[
  \operatorname{spec}L_{\Kc_4}=\{0,\ 1,\ 3,\ 4\}.
\]
\end{proposition}
\begin{proof}
Direct expansion of $\det(L_{\Kc_4}-\lambda I)$ gives
$\lambda^4-8\lambda^3+19\lambda^2-12\lambda
 =\lambda(\lambda-1)(\lambda-3)(\lambda-4)$.
The eigenvalue $0$ is forced: $L_{\Kc_4}\mathbf 1=0$ since every row
sums to zero. Its multiplicity is $1$ because the $1$-skeleton is
connected (Proposition~\ref{prop:b0kernel}).
\end{proof}

The integrality here is a happy accident of this particular graph and
should not be expected in general. It is why $\Kc_4$ was chosen: every
quantity in Parts I--IV can be carried through in exact arithmetic.

% =====================================================================
\chapter{The Sheaf}\label{ch:sheaf}
% =====================================================================

\begin{intu}\Intu
On top of the shape we lay \emph{data}. Each organ holds a vector
describing what it currently believes. Each binding holds a space in
which the beliefs of its members are meant to be compared. A sheaf is
exactly this: local data on every piece, plus rules for moving data into
the places where comparison happens. A ``coherent thought'' is then a
choice of data that agrees everywhere it is compared.
\end{intu}

\section{Which way do the restriction maps point?}\label{sec:direction}

Before the definition, one question must be settled, because the source
documents answer it in two incompatible ways and the answer determines
everything downstream.

Given a face relation $\sigma\triang\tau$ --- say a vertex $u$ on an edge
$e$ --- a sheaf must supply a linear map between $\stalk{u}$ and
$\stalk{e}$. Which direction?

\begin{center}
\begin{tikzpicture}[scale=1.1,
  box/.style={draw,rounded corners=2pt,minimum width=1.9cm,
              minimum height=0.75cm,font=\small},
  lbl/.style={font=\footnotesize}]
  % left: sheaf (correct)
  \node[box] (Fu) at (0,0)    {$\stalk{u}$};
  \node[box] (Fv) at (2.6,0)  {$\stalk{v}$};
  \node[box] (Fe) at (1.3,1.6){$\stalk{e}$};
  \draw[-{Latex[length=2.4mm]},thick,examplecol] (Fu) -- (Fe);
  \draw[-{Latex[length=2.4mm]},thick,examplecol] (Fv) -- (Fe);
  \node[lbl,font=\footnotesize] at (1.3,-0.75) {\textbf{sheaf} (used here)};
  % right: cosheaf (poster)
  \node[box] (Gu) at (6,0)    {$\stalk{u}$};
  \node[box] (Gv) at (8.6,0)  {$\stalk{v}$};
  \node[box] (Ge) at (7.3,1.6){$\stalk{e}$};
  \draw[-{Latex[length=2.4mm]},thick,honestycol] (Ge) -- (Gu);
  \draw[-{Latex[length=2.4mm]},thick,honestycol] (Ge) -- (Gv);
  \node[lbl,font=\footnotesize] at (7.3,-0.75)
       {\textbf{cosheaf} (poster \S4)};
\end{tikzpicture}
\captionof{figure}{The two possible orientations. Only the left one makes
$\ker\cob$ the space of agreeing assignments.}
\label{fig:direction}
\end{center}

\begin{proposition}[The direction is forced]\label{prop:direction}
\Proved\ Suppose we want a linear operator $\cob$ from vertex data to
edge data whose kernel is exactly the set of assignments
$x=(x_v)_{v\in V}$ that ``agree on every edge''. Then the structure maps
must go \emph{from faces to cofaces}, i.e.
$\rest{u\triang e}:\stalk{u}\to\stalk{e}$.
\end{proposition}
\begin{proof}
An assignment gives one vector $x_v\in\stalk{v}$ per vertex, and nothing
on edges. For the two endpoint values $x_u,x_v$ of an edge $e=\{u,v\}$
to be \emph{compared} at all, they must be brought into a common space.
The only space available to $e$ is $\stalk{e}$. Hence we require maps
$\stalk{u}\to\stalk{e}$ and $\stalk{v}\to\stalk{e}$, and agreement on
$e$ means
$\rest{u\triang e}(x_u)=\rest{v\triang e}(x_v)$.

Suppose instead the maps ran downward,
$\rho_{e\to u}:\stalk{e}\to\stalk{u}$. Then from vertex data alone no
map into $\stalk{e}$ exists, so no operator $\Czero\to\Cone$ can be
built from the structure maps, and the condition ``$x$ agrees on $e$''
is not expressible. Downward maps instead assemble into a boundary
operator $\partial:\Cone\to\Czero$, whose kernel consists of \emph{edge}
data summing to zero at vertices --- a flow condition, not an agreement
condition. That is the defining structure of a \emph{cellular cosheaf},
and its invariants are homology $H_\bullet$, not the cohomology
$H^\bullet$ used throughout this book.
\end{proof}

\begin{hon}\Hon
\Refuted\ The poster section \texttt{POSTER/sections/05\_sheaf.tex}
defines, for $\sigma\subseteq\tau$, a map
$\rho_{\tau\to\sigma}:\stalk{\tau}\to\stalk{\sigma}$, and its diagram
draws $\stalk{e}$ above the vertices with arrows pointing down. It then
asserts, three boxes later, that $s\in\Hzero(\CC,\Fs)$ is a global
section. By Proposition~\ref{prop:direction} these two statements are
incompatible: with downward maps the object defined is a cosheaf and
$\Hzero$ as ``the space of compatible assignments'' does not arise. The
main mathematical description
(\texttt{DOCS/TATIANA\_Mathematical\_Description.tex}, Def.~3.1) has the
direction correct. The poster is wrong and should be corrected before it
is shown to anyone. Logged as Appendix~\ref{app:discrepancy}, item D2.
\end{hon}

\begin{exercise}
Take the downward convention seriously for one minute and try to build
$\cob:\Czero\to\Cone$ from it. Where exactly does the construction
fail --- at the definition of $\cob$, or later at the claim
$\ker\cob=\{\text{agreeing assignments}\}$? Answering this is the
fastest way to internalise why the direction is not a convention.
\end{exercise}

\section{The definition}

\begin{definition}[Cellular sheaf]\label{def:sheaf}
A \emph{cellular sheaf} $\Fs$ of real inner-product spaces on a
simplicial complex $\CC$ assigns
\begin{enumerate}[label=(\roman*)]
  \item to each simplex $\sigma\in\CC$ a finite-dimensional real
        inner-product space $\stalk{\sigma}$, the \emph{stalk} over
        $\sigma$;
  \item to each face relation $\sigma\triang\tau$ a linear
        \emph{restriction map}
        \[
          \rest{\sigma\triang\tau}:\stalk{\sigma}\longrightarrow\stalk{\tau},
        \]
\end{enumerate}
subject to functoriality: $\rest{\sigma\triang\sigma}=\Id$, and for every
chain $\sigma\triang\tau\triang\upsilon$,
\begin{equation}\label{eq:functoriality}
  \rest{\tau\triang\upsilon}\circ\rest{\sigma\triang\tau}
  \;=\;\rest{\sigma\triang\upsilon}.
\end{equation}
Equivalently, $\Fs$ is a functor from the face poset of $\CC$, viewed as
a category, to $\Vect_\R$.
\end{definition}

\begin{remark}[Functoriality is a real constraint]\label{rem:functoriality}
Condition~\eqref{eq:functoriality} is vacuous on a complex of dimension
$\le 1$, since there are no chains of length two. It becomes a genuine
constraint the moment a $2$-simplex is filled --- which $\Kc_4$ does. Any
procedure that manufactures restriction maps edge-by-edge (for instance
by fitting each $\rest{v\triang e}$ independently to data) will in
general \emph{violate} \eqref{eq:functoriality} on filled triangles, and
the resulting object is then not a sheaf. This failure mode is not
hypothetical; it is the standard obstruction encountered when
$L^2$-projections are used as restriction maps, and the canonical repair
is to promote the offending interface to a cell of its own so that the
composite is not required to commute. Part~VI returns to
this.
\end{remark}

\begin{definition}[Cochains]\label{def:cochains}
For $k\ge 0$ set
\[
  \Ck{k}\;=\;\bigoplus_{\sigma\in\CC_k}\stalk{\sigma},
\]
the space of \emph{$k$-cochains}, with the direct-sum inner product
$\inner{x}{y}=\sum_\sigma\inner{x_\sigma}{y_\sigma}$. An element of
$\Czero$ is an assignment of one vector to each organ; the current state
$s$ of Definition~\ref{def:M} lives here.
\end{definition}

\section{What the stalks contain}

\begin{construction}[Deployed stalk data]\label{con:stalks}
\Implemented\ In the deployed system every vertex and edge stalk is
$\Rd$ with $d=384$, and the vector $x_v\in\stalk{v}$ is a sentence
embedding produced by a fixed local encoder. Two invariants are enforced
at the level of types:
\begin{enumerate}[label=(\roman*)]
  \item dimensions never truncate or pad --- a mismatch raises rather
        than silently reshaping;
  \item an idle organ carries \emph{no} stalk value, a formal $\unk$,
        never the zero vector.
\end{enumerate}
\end{construction}

\begin{hon}\Hon
Invariant (ii) is the sheaf-theoretic content of ``do not fabricate
data'', and it is worth stating why it is not pedantry. The zero vector
is a genuine point of $\Rd$: it has a distance to every other point, it
contributes to $\norm{X}_F^2$, and it participates in every average.
Substituting $0$ for ``unknown'' therefore does not represent absence,
it asserts a specific belief --- that the organ believes the origin ---
and drags every downstream distance toward it. The distinction between
$\unk$ and $0$ is the difference between a measurement not made and a
measurement made and found to be zero.
\end{hon}

\begin{remark}[Where the fine complex lives]\label{rem:finevsstalk}
The logbook occasionally describes concepts as living \emph{inside an
organ's stalk}. This is loose. The stalk $\stalk{v}$ is a
$384$-dimensional real vector space; the fine complex $\CC_v$ is a
simplicial complex whose vertices are concepts. They are different
objects, and the relation between them is the coarse-graining map
$\pi_v$ of Part~III, which sends the concept data of
$\CC_v$ to a single vector in $\stalk{v}$. Conflating them makes $\pi_v$
look like an identity when it is in fact a lossy precision-weighted
fusion.
\end{remark}

% =====================================================================
\chapter{Coboundary and Laplacian}\label{ch:laplacian}
% =====================================================================

\section{The coboundary}

\begin{definition}[Coboundary in degree zero]\label{def:coboundary}
Fix an arbitrary orientation of each edge, writing $e=(u,v)$ for the
edge $\{u,v\}$ oriented from $u$ to $v$. The \emph{coboundary}
$\cob:\Czero\to\Cone$ acts edgewise by
\begin{equation}\label{eq:coboundary}
  (\cob x)_e\;=\;\rest{v\triang e}(x_v)\;-\;\rest{u\triang e}(x_u)
  \ \in\ \stalk{e}.
\end{equation}
A $0$-cochain $x$ is a \emph{global section} if $\cob x=0$, that is, if
every bound pair agrees after being mapped into its shared comparison
space.
\end{definition}

\begin{lemma}[Orientation-independence]\label{lem:orientation}
\Proved\ Reversing the orientation of an edge $e$ replaces
$(\cob x)_e$ by its negative. Consequently $\ker\cob$, the quantity
$\norm{\cob x}^2$, and the operator $\cob^\top\cob$ are all independent
of the chosen orientations.
\end{lemma}
\begin{proof}
Reversal exchanges the roles of $u$ and $v$ in \eqref{eq:coboundary},
negating the entry. Negating a coordinate leaves its vanishing, its
squared norm, and hence the sum of squared norms unchanged. Writing
$\cob'=D\cob$ with $D$ diagonal with entries $\pm1$ blockwise, we get
$\cob'^\top\cob'=\cob^\top D^\top D\cob=\cob^\top\cob$ since $D^\top D=I$.
\end{proof}

\begin{remark}[Why an arbitrary choice is harmless]
Lemma~\ref{lem:orientation} is what licenses the phrase ``fix an
arbitrary orientation''. It is worth checking rather than assuming,
because the analogous statement fails for the \emph{signed} per-edge
quantity $(\cob x)_e$ itself, which is only defined up to sign. Any
diagnostic that reports $(\cob x)_e$ rather than $\norm{(\cob x)_e}^2$
is reporting an orientation-dependent quantity and is therefore not
well defined.
\end{remark}

\section{The sheaf Laplacian}

\begin{definition}[Sheaf Laplacian and discord]\label{def:laplacian}
The \emph{degree-zero sheaf Laplacian} is
\begin{equation}\label{eq:laplacian}
  \Lap\;=\;\cob^\top\cob\;:\;\Czero\to\Czero,
\end{equation}
and the \emph{discord} of a state $x$ is the Dirichlet energy
\begin{equation}\label{eq:omega}
  \discord(x)\;=\;x^\top \Lap x\;=\;\norm{\cob x}^2
  \;=\;\sum_{e=\{u,v\}}
       \norm{\rest{v\triang e}(x_v)-\rest{u\triang e}(x_u)}^2 .
\end{equation}
\end{definition}

\begin{proposition}[Basic properties]\label{prop:Lpsd}
\Proved\ $\Lap$ is symmetric positive semidefinite, and
\[
  \discord(x)=0
  \iff \cob x=0
  \iff x\in\ker \Lap .
\]
Moreover $\ker\cob=\Hzero$, the zeroth sheaf cohomology.
\end{proposition}
\begin{proof}
Symmetry: $(\cob^\top\cob)^\top=\cob^\top\cob$. Positive
semidefiniteness: $x^\top \Lap x=\norm{\cob x}^2\ge0$. If $\Lap x=0$
then $\norm{\cob x}^2=x^\top \Lap x=0$, so $\cob x=0$ since
$\norm{\cdot}$ is a norm; conversely $\cob x=0$ gives $\Lap x=0$
immediately. Finally the cochain complex begins in degree $0$ --- there
are no $(-1)$-cochains --- so $\Hzero=\ker\cob^0/\im\cob^{-1}=\ker\cob$.
\end{proof}

\begin{intu}\Intu
$\discord$ is a single number: the total squared disagreement across all
bindings. It vanishes exactly when the society holds one coherent
thought. The point of the construction is that ``a coherent thought'' has
become a precise object --- a vector in the kernel of a matrix --- rather
than a mood.
\end{intu}

\begin{proposition}[Kernel and components]\label{prop:b0kernel}
\Proved\ If every restriction map is the identity on a common stalk
$\Rd$, then
\[
  \Lap=L_{\Kc}\otimes I_d,
\]
where $L_{\Kc}$ is the ordinary combinatorial graph Laplacian of the
$1$-skeleton. Consequently $\operatorname{spec}\Lap$ is
$\operatorname{spec}L_\Kc$ with each eigenvalue repeated $d$ times, and
$\dim\ker\Lap=d\cdot b_0$, where $b_0$ is the number of connected
components.
\end{proposition}
\begin{proof}
With identity restrictions \eqref{eq:omega} reads
$\discord(x)=\sum_{\{u,v\}\in E}\norm{x_u-x_v}^2$, which is the
Dirichlet form of $L_\Kc$ applied coordinatewise in each of the $d$
components. Hence $\Lap=L_\Kc\otimes I_d$. The spectrum of a Kronecker
product with the identity is the spectrum of the factor with multiplicity
$d$. The kernel of $L_\Kc$ is spanned by the indicator vectors of the
connected components, of which there are $b_0$; tensoring with $\Rd$
multiplies the dimension by $d$.
\end{proof}

\begin{hon}\Hon
\Implemented\ The deployed system uses identity restriction maps, so
Proposition~\ref{prop:b0kernel} describes the regime it actually
operates in. This is honest and cheap, but it is the \emph{crudest}
possible sheaf: on a connected $1$-skeleton, $\ker\Lap$ is the diagonal
$\{x:x_u=x_v\ \forall u,v\}$, so the only coherent states are exact
consensus and reconciliation is plain averaging. Everything interesting
that a sheaf can express --- that two organs agree \emph{after} a
change of basis, or agree only on a subspace --- lives in non-identity
restriction maps. Promoting them (learned $\rest{v\triang e}$, for
instance by closed-form Procrustes alignment) is the identified route
from a diagnostic sheaf to a computational one. It is a direction, not a
result.
\end{hon}

\section{The running example, computed}

\begin{workedbox}\Worked
Take $\Kc_4$ from Example~\ref{ex:K4} with identity restrictions and
stalk dimension $d=2$, and give it the state
\[
  x_{\mathsf S}=\begin{pmatrix}1\\0\end{pmatrix},\quad
  x_{\mathsf M}=\begin{pmatrix}1\\0\end{pmatrix},\quad
  x_{\mathsf R}=\begin{pmatrix}0\\1\end{pmatrix},\quad
  x_{\mathsf V}=\begin{pmatrix}0\\1\end{pmatrix}.
\]
Read it as: \textsc{Search} and \textsc{Memory} agree with each other,
\textsc{Reasoning} and \textsc{Verify} agree with each other, and the two
camps are orthogonal.

\medskip
\emph{Discord, edge by edge.} Using \eqref{eq:omega} with identity maps,
$\discord_e=\norm{x_u-x_v}^2$:
\[
\begin{array}{lcl}
  \{\mathsf S,\mathsf M\}:&\ \norm{(0,0)}^2 &= 0,\\
  \{\mathsf M,\mathsf R\}:&\ \norm{(1,-1)}^2 &= 2,\\
  \{\mathsf S,\mathsf R\}:&\ \norm{(1,-1)}^2 &= 2,\\
  \{\mathsf R,\mathsf V\}:&\ \norm{(0,0)}^2 &= 0,
\end{array}
\qquad\Longrightarrow\qquad \discord(x)=4 .
\]

\emph{Localisation.} The two guilty edges tie at $\discord_e=2$. The
$\arg\max$ is therefore \emph{not unique}, and a diagnostic that returns
a single ``worst edge'' here is returning an artefact of tie-breaking
order, not a measurement. The honest return is the tied set
$\{\{\mathsf M,\mathsf R\},\{\mathsf S,\mathsf R\}\}$ --- which, read
back, correctly says that $\mathsf R$ is the organ at odds with the rest.

\emph{Kernel projection.} The $1$-skeleton is connected, so by
Proposition~\ref{prop:b0kernel} the kernel is the diagonal and $x_0$ is
the constant assignment at the mean
$\bar x=\tfrac14\bigl((1,0)+(1,0)+(0,1)+(0,1)\bigr)
 =\bigl(\tfrac12,\tfrac12\bigr)$.
Hence
\[
  \norm{x_0}^2=4\cdot\Bigl(\tfrac14+\tfrac14\Bigr)=2,
  \qquad
  x_\perp = x-x_0
  =\Bigl(\pm\tfrac12,\mp\tfrac12\Bigr)_{v},
  \qquad
  \norm{x_\perp}^2=4\cdot\tfrac12=2,
\]
and Pythagoras checks: $\norm{X}_F^2=2+2=4$. \hfill$\square$
\end{workedbox}

These numbers are picked up again in Part~II, where they are fed to the
coherence score $\rhoscore$ and its factorisation, and in
Part~IV, where the same complex is used to compute
$b_0$, $b_1$ and $\Hone$.

\begin{exercise}
Recompute $\discord(x)$ for the same complex after filling in a
restriction map that is \emph{not} the identity: let
$\rest{\mathsf R\triang e}$ be rotation by $90^\circ$ for the two edges
at $\mathsf R$ inside the triangle, and the identity elsewhere. Does
$\discord$ go up or down, and does the resulting assignment of maps
satisfy functoriality \eqref{eq:functoriality} on the filled triangle
$\{\mathsf S,\mathsf M,\mathsf R\}$? The second half of that question is
the important one.
\end{exercise}
